\documentclass[11pt,a4paper]{article}
\usepackage[utf8]{inputenc}
\usepackage{hyperref}
\usepackage{geometry}
\geometry{margin=1in}

\title{Glance ML Internship Assignment: \\ Multimodal Fashion \& Context Retrieval}
\author{}
\date{}

\begin{document}

\maketitle

\section{Approaches}
To build a system capable of searching fashion images by specific clothing attributes, environment context, and style vibe, several approaches were evaluated:

\subsection*{A. Zero-Shot Vision-Language Models (e.g., CLIP, SigLIP)}
\begin{itemize}
    \item \textbf{How it works:} Encodes both images and search queries into a shared multi-modal embedding space. Retrieval is performed using cosine similarity.
    \item \textbf{Trade-offs:} Highly scalable and fast. However, standard CLIP models struggle severely with compositionality (e.g., distinguishing ``red shirt, blue pants'' from ``blue shirt, red pants'') and often ignore fine-grained fashion terminology.
    \item \textbf{When to use:} Good as a baseline for broad semantic searches, but insufficient for strict attribute constraints.
\end{itemize}

\subsection*{B. Object Detection + Attribute Classification (e.g., Faster R-CNN)}
\begin{itemize}
    \item \textbf{How it works:} Detects bounding boxes for clothing items and runs a classifier to predict color, fabric, and style for each box.
    \item \textbf{Trade-offs:} Highly accurate and compositional. However, the pipeline is extremely rigid, relies heavily on predefined bounding-box datasets, and scales poorly to open-vocabulary natural language queries.
    \item \textbf{When to use:} Ideal for structured e-commerce databases, but fails for contextual ``vibe'' or ``environment'' searches.
\end{itemize}

\subsection*{C. Hybrid Semantic Retrieval with Fine-Tuned VLM (The Chosen Approach)}
\begin{itemize}
    \item \textbf{How it works:} Fuses the strengths of state-of-the-art vision encoders (FashionSigLIP) with Large Vision-Language Models (Qwen3-VL) that have been explicitly fine-tuned to extract both visual embeddings and highly detailed, ontology-aware text captions for each image.
    \item \textbf{Trade-offs:} Requires significant upfront compute during the training and indexing phases, but provides unparalleled retrieval accuracy for compositional, attribute-heavy, and context-aware queries.
    \item \textbf{When to use:} Best for complex, multi-attribute, open-vocabulary semantic search.
\end{itemize}

\section{Short Write-up on Chosen Approach}

The implemented system leverages a \textbf{Hybrid Semantic Retrieval Pipeline}, supported by a massive custom training orchestration.

\subsection*{1. Model Fine-Tuning \& Evaluation}
Standard VLMs lack the strict vocabulary required for fashion retrieval. To solve this, we orchestrated a distributed Supervised Fine-Tuning (SFT) pipeline targeting \texttt{Qwen3-VL-8B-Thinking}:
\begin{itemize}
    \item \textbf{Dataset Processing:} A dataset of 8,500 images was processed and sliced into customized subsets (5,000 for SFT, 1,000 for Validation/Test). Image inputs were formatted using HuggingFace Chat Templates.
    \item \textbf{Cloud Execution (Modal):} To handle the heavy compute, the training pipeline was deployed to a serverless \textbf{A100 (40GB) GPU} instance using Modal.
    \item \textbf{LoRA Integration:} We applied Low-Rank Adaptation (\texttt{r=16}) to all dense layers (\texttt{q\_proj, k\_proj, v\_proj, o\_proj, gate\_proj, up\_proj, down\_proj}) to efficiently train the 8B model. 
    \item \textbf{Custom Collator \& M-RoPE:} We built a custom data collator to dynamically patch visual matrices (\texttt{image\_grid\_thw}) and inject \texttt{mm\_token\_type\_ids} required by Qwen3's Multimodal Rotary Position Embedding (M-RoPE) mechanics.
    \item \textbf{Evaluation Metrics:} The SFT model was evaluated on the 1,000-image test set using an exacting substring-match algorithm. It achieved an \textbf{Attribute F1 Score of 0.0791 (7.9\%)}. While exact-match string evaluations are notoriously strict for generative outputs, this provided a strong empirical baseline confirming the model learned the prompt structures. (BERTScore integration was also added to track semantic improvements).
\end{itemize}

\subsection*{2. The Indexing Architecture}
For every image in the database, the indexer extracts two representations:
\begin{enumerate}
    \item \textbf{Visual Embeddings:} We use \texttt{Marqo-FashionSigLIP}, a variant of SigLIP explicitly fine-tuned on fashion datasets, to generate a 768-dimensional visual embedding. 
    \item \textbf{Semantic Captions:} The fine-tuned \texttt{Qwen3-VL} generates an incredibly detailed structured caption. It explicitly identifies environments, vibes, clothing patterns, colors, and accessories. These captions are subsequently embedded into text vectors.
\end{enumerate}
Both embeddings, along with the raw metadata, are stored in \textbf{ChromaDB}.

\subsection*{3. The Retrieval Architecture}
When a user issues a natural language query (e.g., ``A red tie and a white shirt in a formal setting''):
\begin{enumerate}
    \item The query is embedded using the FashionSigLIP text encoder.
    \item The engine performs an approximate nearest neighbor search across both the visual embedding collection and the VLM caption embedding collection.
    \item A \textbf{Hybrid Scoring Function} is applied: \\
    \texttt{Score = ($\alpha$ $\times$ Visual Similarity) + ($\beta$ $\times$ Caption Similarity) + ($\gamma$ $\times$ Metadata Boost)}
    \item The Metadata Boost heavily penalizes candidates that lack the specific colors or objects explicitly mentioned in the query, effectively solving standard CLIP compositionality failures.
\end{enumerate}

\section{Codebase Link}
The complete codebase is documented and structured cleanly. \\
\url{https://github.com/vidit68/glance-multimodal-retrieval} \\
\textit{Note: The codebase heavily utilizes Modal serverless infrastructure to orchestrate massive VLM fine-tuning across A100 GPUs and handles data via high-volume persistent mounts.}

\section{Approaches for Future Work}

\subsection*{A. Implementing Reinforcement Learning (GSPO) and Custom Loss}
\begin{itemize}
    \item \textbf{RL Alignment via GSPO:} While we established the SFT baseline, future work involves completing the Group Sequence Policy Optimization (GSPO) pipeline. By employing a reward function that directly evaluates the generation quality (F1 / BERTScore), we can align the VLM strictly to fashion retrieval metrics.
    \item \textbf{Custom Loss Formulation:} We intend to design a custom loss function during the GSPO phase that dynamically weights precision over recall, heavily penalizing hallucinated attributes (e.g., predicting ``red tie'' when none exists) compared to missed minor attributes. This custom loss will seamlessly integrate with the policy gradient updates.
\end{itemize}

\subsection*{B. Adding Locations (Cities, Places) and Weather}
\begin{itemize}
    \item \textbf{Geospatial \& Temporal Indexing:} If the dataset contains GPS coordinates or timestamps, we can integrate a spatial-temporal vector database.
    \item \textbf{VLM Prompt Expansion:} We can explicitly prompt the indexing VLM to infer weather (e.g., ``sunny'', ``winter snow'') and architectural location cues from the image background. 
\end{itemize}

\subsection*{C. Improving Precision with Cross-Encoders}
\begin{itemize}
    \item \textbf{Cross-Encoder Re-Ranking:} While the current system uses bi-encoders (SigLIP) for fast retrieval, precision can be significantly improved by applying a heavy Cross-Encoder (like ALBEF or BLIP-2 ITM) to the top-50 results. The cross-encoder can attend to the interaction between query words and image patches, entirely eliminating compositionality errors.
\end{itemize}

\end{document}
